\section{Conclusion}

The identification of functional protein modules in PPI networks remains challenging due to the sparsity and incompleteness of interaction data, as well as the overlapping and nested organization of protein complexes. Addressing these challenges requires methods that effectively integrate the structural information encoded in PPI networks with biological knowledge associated with individual proteins. In this study, we proposed an end-to-end learning framework for functional module identification that combines GNNs with GO information. In particular, we investigated the complementary use of binary GO features and GO2Vec embeddings. This strategy improved the overall performance compared with using either feature representation independently.

Furthermore, we extended the NOCD loss function in two directions to better account for the characteristics of PPI networks. The first extension, HO-NOCD, incorporates higher-order connectivity information, while the second, W-NOCD, incorporates interaction weights to account for the varying strengths of protein interactions. Experimental comparisons demonstrated that W-NOCD achieved the best overall performance among the investigated loss functions, highlighting the benefit of incorporating interaction strength into the learning objective. We also investigated different GNN architectures, aggregation mechanisms, and activation functions. The results indicated that a standard two-layer GAT architecture with ReLU activation, without additional JKNet aggregation mechanisms, provided the most effective configuration for this task.

To address the overlapping and nested organization of protein complexes, we further introduced an aggregated thresholding mechanism that combines module predictions obtained using multiple affiliation thresholds. This strategy improved the coverage of gold-standard complexes, particularly for nested module pairs, demonstrating its potential to recover relationships between complexes that may be missed when using a single fixed threshold. In addition, combining binary GO features with GO2Vec embeddings improved the overall F1-score and recall, despite a reduction in precision. This finding suggests that the two feature representations provide complementary functional information, enabling the identification of additional gold-standard complexes.

Finally, we compared the proposed framework with existing methods across five yeast PPI datasets. The results demonstrated that our approach consistently achieved higher F1-scores than the competing methods, indicating the effectiveness of jointly exploiting network topology, interaction weights, and protein functional annotations for functional module identification.

Despite these promising results, the proposed framework has primarily been evaluated on yeast PPI networks, consistent with the scope of many existing studies. An important direction for future research is to investigate its generalizability to human PPI networks and extend its application to related biological tasks, such as pathway identification. Another direction concerns the representation of GO information in GO2Vec. In this study, mean pooling was applied to the embeddings of each GO subspace, followed by concatenation to obtain the final representation. Although this strategy is computationally efficient, mean pooling may discard potentially informative relationships among GO terms. Therefore, exploring alternative pooling strategies and more expressive aggregation mechanisms could further improve the representation of functional information. Further research into more effective ways of integrating information across GO subspaces may provide additional improvements.